\section*{Schedules}
\label{sec:schedules}

\scheduleSection{Introduction to number systems and logic}{2}{
Overview of the natural numbers, integers, real numbers, rational and irrational numbers, algebraic and transcendental numbers. Brief discussion of complex numbers; statement of the Fundamental Theorem of Algebra.

\noindent Ideas of axiomatic systems and proof within mathematics; the need for proof; the role of counter-examples in mathematics. Elementary logic; implication and negation; examples of negation of compound statements. Proof by contradiction.}

\scheduleSection{Sets, relations and functions}{4}{
Union, intersection and equality of sets. Indicator (characteristic) functions; their use in establishing set identities. Functions; injections, surjections and bijections. Relations, and equivalence relations. Counting the combinations or permutations of a set. The Inclusion-Exclusion Principle.}

\scheduleSection{The integers}{2}{
The natural numbers: mathematical induction and the well-ordering principle. Examples, including the Binomial Theorem.}

\scheduleSection{Elementary number theory}{8}{
Prime numbers: existence and uniqueness of prime factorisation into primes; highest common factors and least common multiples. Euclid's proof of the infinity of primes. Euclid's algorithm. Solution in integers of \(ax + by = c\).

\noindent Modular arithmetic (congruences). Units modulo \(n\). Chinese Remainder Theorem. Wilson's Theorem; the Fermat-Euler Theorem. Public key cryptography and the RSA algorithm.}

\scheduleSection{The real numbers}{4}{
Least upper bounds; simple examples. Least upper bound axiom. Sequences and series; convergence of bounded monotonic sequences. Irrationality of \(\sqrt{2}\) and \(e\). Decimal expansions. Construction of a transcendental number.}

\scheduleSection{Countability and uncountability}{4}{
Definitions of finite, infinite, countable and uncountable sets. A countable union of countable sets is countable. Uncountability of \(\RR\). Non-existence of a bijection from a set to its power set. Indirect proof of existence of transcendental numbers.}